\section{Who Owns a Field}
\label{sec:ch05-who-owns-a-field}

\index{ownership!the rule|(}%
The rule is one sentence. A field belongs to the subgraph that owns the data
behind it, and no other subgraph declares it.

\index{ownership!an edge and its target are owned separately}%
Apply it to \code{Session.speaker} and it splits in a way that is easy to get
backwards. The field is on \code{Session}, and which speaker gives a session is
a fact the Sessions service stores, in a column, next to the row. So Sessions
owns \code{speaker}. What Sessions does not own is the speaker: the name, the
biography, everything the field's type promises. Sessions declares the field,
produces the key, and stops there. The edge and its target have different
owners, and the seam runs between them rather than through either.

\index{ownership!a subgraph contributing to a type it does not own}%
The second arrangement is the one that makes federation more than remote
procedure calls. The Ratings service stores nobody's sessions and nobody's
speakers, and it is nevertheless the right place for a session's ratings to
live. So it declares \code{Session} too, carrying the key and two of the three
fields chapter~\ref{ch:second-third-subgraph} gives it; the third needs a
directive section~\ref{sec:ch05-the-directives} has not introduced yet:

\begin{graphqlsdl}
type Session @key(fields: "id") {
  id: ID!
  averageScore: Float
  ratingCount: Int!
}
\end{graphqlsdl}

\index{Session@\code{Session}!assembled from two services}%
Compose that with the Sessions subgraph and a client sees one \code{Session}
carrying eight fields, two of which come from a service that stores no
sessions, and nothing in the schema says which came from where. A client asking
for a title and an average score has asked one question of one endpoint. That
is what the architecture is for, and it is delivered by a type appearing twice
with different fields on it.

\index{ownership!is a claim the composer enforces}%
Which is why ownership is not a convention here. It is a claim each subgraph
makes in its schema, and the composer checks the claims against each other
before anything runs. Chapter~\ref{ch:do-you-need-this} argued that federation
does not remove coordination between teams but moves it into what a program
checks. This is that program, and these are the checks.

\subsection{When Two Subgraphs Claim the Same Field}

\index{shareable@\code{"@shareable}!declares a field may be duplicated}%
Declare a field in two subgraphs and composition fails. Not with a warning, not
with a last-one-wins: it fails, and the message names the field and every
subgraph that declared it, and says that none of them declared it
\code{@shareable}. The composer's position is that a field defined twice is
either a mistake or a decision, and it declines to guess which.

\index{shareable@\code{"@shareable}!is a promise about the value}%
\code{@shareable} is how you say it was a decision. It means the field may be
resolved by more than one subgraph and that each of them will return the same
value for the same object, which is a promise about your data and not a
compiler flag. Where that promise holds, the router gains a choice: it can get
the field from whichever subgraph it was already talking to and save a hop.
Where it does not hold, you have built a graph that returns different answers
depending on a routing decision no client can see.

\index{node@\code{node}!collides across subgraphs}%
This book has already bought itself one of these, and the place to meet it is
here rather than in the middle of chapter~\ref{ch:first-subgraph}.
Chapter~\ref{ch:schema-design} turned on global object identification, which
puts \code{node} and \code{nodes} on \code{Query}. Every subgraph in this book
will therefore export both. Two subgraphs, two \code{Query.node} fields, and
composition stops with exactly the error above. The fix is to declare them
shareable in both, and the interesting part is that the promise is real: any
subgraph handed a global node id has to answer with the same object, which is
the promise chapter~\ref{ch:schema-design}'s id format already makes.

\subsection{Moving a Field Between Owners}

\index{override@\code{"@override}!moves a field}%
Ownership changes. A field that started in the service that had the data ends
up in the service that grew up around it, and the change has to happen without
a coordinated deployment, because avoiding coordinated deployments is why you
federated. \code{@override(from:)} is the mechanism: the receiving subgraph
declares the field and names the subgraph it is taking it from, composition
gives it the field, and the old subgraph keeps its copy until somebody deletes
it.

\index{override@\code{"@override}!two deployments, either order}%
That ordering is the point. Both subgraphs carry the field for as long as you
like, so the two deployments happen on two teams' calendars in either order,
and the graph is answerable throughout. It is the same problem
chapter~\ref{ch:schema-design} solved for a field a client depends on, one level
up, with services in the place of clients.

\index{ownership!the part no directive settles}%
None of which decides who should own a field, and that is not an oversight in
the specification. \code{@shareable} lets two teams both be right;
\code{@override} lets one take a field from another; neither has an opinion
about whether they should. Chapter~\ref{ch:ownership} is the chapter about the
part that stays a conversation, and I put it in part~III with the failures
rather than here with the mechanisms because that is where I think it belongs.
\index{ownership!the rule|)}%
